%=====================================================================
\chapter{The Data Science Lifecycle}\label{ch:lifecycle}
%=====================================================================
\locator{1}{Part I --- Foundations of Data}

\begin{outcomes}
By the end of this chapter you will be able to:
\begin{tight}
  \item \textbf{Describe} the six phases of the CRISP-DM lifecycle and explain why the process is a loop rather than a line.
  \item \textbf{Translate} a vague business request into a precise, measurable data science problem statement.
  \item \textbf{Distinguish} a business metric from a model metric, and explain why a project needs both.
  \item \textbf{Estimate} where project effort actually goes, and plan accordingly.
  \item \textbf{Recognise} the four failure modes that kill data science projects before any model is fitted.
\end{tight}
\end{outcomes}

\begin{prereq}
\chref{landscape} --- particularly the four levels of analytics, which determine
which lifecycle phases a project will need.
\end{prereq}

\begin{keyterms}
CRISP-DM $\bullet$ problem framing $\bullet$ business understanding $\bullet$ data
understanding $\bullet$ baseline $\bullet$ business metric $\bullet$ model metric
$\bullet$ success criterion $\bullet$ deployment $\bullet$ iteration
\end{keyterms}

\section{Why a Lifecycle at All}

Beginners imagine data science as: get data, build model, done. Practitioners
know it as a loop that revisits its own assumptions repeatedly, and in which the
modelling step is usually the shortest.

The value of naming the phases is not bureaucracy. It is that each phase has a
distinct \emph{failure mode}, and being able to say ``we are stuck in data
understanding'' is the first step to getting unstuck.

\section{CRISP-DM: The Standard Lifecycle}

\begin{definitionbox}[CRISP-DM]
The \textbf{Cross-Industry Standard Process for Data Mining}: a six-phase,
explicitly cyclic process model --- business understanding, data understanding,
data preparation, modelling, evaluation and deployment. It remains the most widely
taught lifecycle because it puts \emph{understanding the problem} before
\emph{touching the data}.
\end{definitionbox}

\dsfig{%
\begin{tikzpicture}[scale=1.0,
  ph/.style={circle, draw=#1, line width=1.2pt, fill=#1!10, text=#1!50!black,
             align=center, font=\scriptsize\bfseries, minimum size=2.15cm, inner sep=1pt},
  arr/.style={-{Stealth[length=2.6mm]}, primary!55, line width=1.1pt}]

  \node[ph=primary]   (bu) at (90:3.3)  {Business\\Understanding};
  \node[ph=secondary] (du) at (30:3.3)  {Data\\Understanding};
  \node[ph=tealhl]    (dp) at (-30:3.3) {Data\\Preparation};
  \node[ph=greenhl]   (mo) at (-90:3.3) {Modelling};
  \node[ph=goldhl]    (ev) at (-150:3.3){Evaluation};
  \node[ph=purplehl]  (de) at (150:3.3) {Deployment};

  \draw[arr] (bu) to[bend left=14] (du);
  \draw[arr] (du) to[bend left=14] (dp);
  \draw[arr] (dp) to[bend left=14] (mo);
  \draw[arr] (mo) to[bend left=14] (ev);
  \draw[arr] (ev) to[bend left=14] (de);
  \draw[arr] (de) to[bend left=14] (bu);

  % the loops that beginners do not expect
  \draw[{Stealth[length=2mm]}-{Stealth[length=2mm]}, accent!70, dashed, line width=0.9pt]
        (bu) -- (du);
  \draw[{Stealth[length=2mm]}-{Stealth[length=2mm]}, accent!70, dashed, line width=0.9pt]
        (dp) -- (mo);
  \draw[{Stealth[length=2mm]}-{Stealth[length=2mm]}, accent!70, dashed, line width=0.9pt]
        (ev) -- (bu);

  \node[font=\scriptsize, text=primary!70!black, align=center] at (0,0.35) {THE DATA};
  \node[font=\scriptsize, text=primary!70!black, align=center] at (0,-0.05) {\bfseries CRISP-DM};
  \node[font=\tiny\itshape, text=accent!80!black, align=center] at (0,-0.55)
       {dashed = the loops\\beginners do not expect};
\end{tikzpicture}}%
{The CRISP-DM lifecycle. The solid arrows are the nominal order; the dashed arrows are the
back-steps that happen in every real project. Evaluation sending you back to business
understanding is not a sign of failure --- it is the process working.}{crispdm}

\subsection{What Happens in Each Phase}

\rowcolors{2}{white}{lightbg}
\begin{longtable}{@{}L{2.9cm}L{5.3cm}L{6.2cm}@{}}
\toprule
\rowcolor{primary!15}
\textbf{Phase} & \textbf{The work} & \textbf{Its characteristic failure} \\
\midrule
\endhead
Business Understanding & Agree the objective, the decision it supports, the success criterion and the constraints & Building something nobody asked for, beautifully \\
Data Understanding & Collect, describe, explore and check the quality of what is available & Discovering in month three that the key field is 60\% empty \\
Data Preparation & Clean, join, transform, engineer features, split the data & Leaking information from the future into the training set \\
Modelling & Select algorithms, train, tune hyperparameters & Chasing a 0.3\% gain that the business cannot perceive \\
Evaluation & Check against the \emph{business} criterion, not just the metric & Declaring victory on a test set that does not resemble reality \\
Deployment & Ship it, monitor it, retrain it, document it & Handing over a notebook and calling it a system \\
\bottomrule
\end{longtable}

\begin{conceptbox}[Why Deployment Points Back to Business Understanding]
The arrow from deployment back to the start is the most important one in
\dsref{crispdm}. Once a model is live it changes the behaviour it was trained to
predict --- students who are warned study differently, attackers who are blocked
change tactics. The world you modelled no longer exists, so the cycle restarts.
\chref{mlops} is entirely about living with this.
\end{conceptbox}

\section{Framing the Problem}

Most projects that fail, fail here. A request arrives in the language of the
business --- ``reduce dropouts'', ``stop the breaches'', ``ship on time'' --- and
must be converted into something a computer can be scored against.

\begin{definitionbox}[Problem Framing]
The act of converting a stakeholder's goal into a precise specification: what is
being predicted or described, for which unit of analysis, using what data
available at what time, judged by what criterion, to support which decision.
\end{definitionbox}

A complete frame answers six questions. If you cannot answer all six, you are not
ready to open the data.

\dsfig{%
\begin{tikzpicture}[
  q/.style={rounded corners=3pt, draw=#1, line width=1pt, fill=#1!7,
            text width=3.5cm, align=center, font=\scriptsize, inner sep=4pt},
  hd/.style={font=\bfseries\scriptsize}]

  \node[q=primary]   (a) at (0,0)      {\textbf{1. DECISION}\\What action will this change? If none: stop.};
  \node[q=secondary] (b) at (4.3,0)    {\textbf{2. UNIT}\\One row is one \emph{what}? Student? Session? Sensor-hour?};
  \node[q=tealhl]    (c) at (8.6,0)    {\textbf{3. TARGET}\\Exactly which quantity is predicted or measured?};
  \node[q=greenhl]   (d) at (0,-2.6)   {\textbf{4. TIMING}\\What is knowable \emph{at prediction time}? Nothing later.};
  \node[q=goldhl]    (e) at (4.3,-2.6) {\textbf{5. CRITERION}\\What number, at what level, means success?};
  \node[q=purplehl]  (f) at (8.6,-2.6) {\textbf{6. BASELINE}\\What does the current, model-free approach achieve?};

  \draw[-{Stealth[length=2mm]}, gray!50] (a) -- (b);
  \draw[-{Stealth[length=2mm]}, gray!50] (b) -- (c);
  \draw[-{Stealth[length=2mm]}, gray!50] (c.south) .. controls +(0,-0.8) and +(0,0.8) .. (d.north);
  \draw[-{Stealth[length=2mm]}, gray!50] (d) -- (e);
  \draw[-{Stealth[length=2mm]}, gray!50] (e) -- (f);
\end{tikzpicture}}%
{The six framing questions. Question~4 (timing) prevents the single most expensive mistake
in applied machine learning --- training on information that will not exist when the model
is actually used. Question~6 stops you celebrating a model that is worse than a rule of
thumb.}{framing-six}

\begin{alertbox}[Question 4 Deserves Special Fear]
A model that predicts student failure using \emph{final attendance percentage}
will look superb in testing and be useless in week~3, because final attendance is
not known in week~3. This is called \term{leakage}, and it is treated in detail in
\chref{wrangling} and \chref{evaluation}. It has ruined more student projects than
any other single error.
\end{alertbox}

\section{Two Kinds of Metric}

\begin{definitionbox}[Business Metric vs.\ Model Metric]
A \term{model metric} scores the model against the data (accuracy, error, AUC). A
\term{business metric} scores the project against the world (dropouts prevented,
downtime avoided, breaches caught, cost saved). A project needs both, and they can
move in opposite directions.
\end{definitionbox}

\begin{worked}[When a Better Model Is a Worse Project]
A dropout-prediction model is improved from 88\% to 91\% accuracy. Is the project
better off?

\smallskip
\textbf{The data.} 1{,}000 students, 100 of whom drop out. The department can fund
50 interventions.

\textbf{Model A} (88\% accurate) flags 60 students, of whom 40 truly drop out.\\
\textbf{Model B} (91\% accurate) flags 20 students, of whom 18 truly drop out.

\textbf{Model metric.} B wins on accuracy: it makes fewer mistakes overall,
largely by predicting ``will not drop out'' more often --- which is right 90\% of
the time by default.

\textbf{Business metric.} The department can afford 50 interventions and each
successful intervention retains one student. Model~A surfaces 40 real cases;
Model~B surfaces 18. Under capacity of 50, Model~A retains up to 40 students,
Model~B up to 18.

\textbf{Conclusion.} Model~A is worth roughly twice as much despite being the
worse model by accuracy. The mismatch arises because accuracy ignores both the
class imbalance and the capacity constraint. \chref{evaluation} builds the metrics
that would have caught this --- here, \emph{recall} at a fixed budget.
\end{worked}

\section{Where the Time Actually Goes}

\dsfig{%
\begin{tikzpicture}
\begin{axis}[
  width=12.2cm, height=5.0cm,
  xbar stacked, xmin=0, xmax=100,
  bar width=15pt,
  ytick={0,1}, yticklabels={{What students\\expect},{What projects\\actually need}},
  yticklabel style={align=center, font=\scriptsize},
  xlabel={\scriptsize share of project effort (\%)},
  xtick={0,20,40,60,80,100},
  tick label style={font=\scriptsize},
  legend style={at={(0.5,-0.42)}, anchor=north, legend columns=5,
                font=\scriptsize, draw=gray!40},
  axis lines=left, enlarge y limits=0.55,
]
\addplot[fill=primary!70, draw=none]   coordinates {(10,0) (25,1)};
\addplot[fill=secondary!65, draw=none] coordinates {(10,0) (35,1)};
\addplot[fill=greenhl!65, draw=none]   coordinates {(60,0) (15,1)};
\addplot[fill=goldhl!70, draw=none]    coordinates {(10,0) (10,1)};
\addplot[fill=purplehl!60, draw=none]  coordinates {(10,0) (15,1)};
\legend{Framing, Data prep, Modelling, Evaluation, Deployment}
\end{axis}
\end{tikzpicture}}%
{Expected versus actual effort. Modelling --- the part that attracted you to the subject
--- is typically the smallest phase. Framing and data preparation together consume most
of a real project, which is why Chapters 2, 3 and 7 exist.}{effort-split}

\begin{pitfall}[The Four Project Killers]
\begin{tight}
  \item \textbf{No decision attached.} If no action changes based on the output,
        the project cannot succeed, however accurate it is.
  \item \textbf{No baseline.} Without knowing what a simple rule achieves, you
        cannot show your model is worth its complexity.
  \item \textbf{Unavailable-at-prediction-time features.} See the alert above.
  \item \textbf{A success criterion agreed only after seeing the results.} If the
        target moves to wherever the model landed, the evaluation is theatre.
\end{tight}
\end{pitfall}

%---------------------------------------------------------------------
\begin{fourlenses}{Problem Framing}
\lensLA{\emph{Decision:} which students receive an outreach call in week 5.
\emph{Unit:} one enrolled student in one course-offering. \emph{Target:} will the
student withdraw or fail? \emph{Timing:} only data from weeks 1--4 may be used.
\emph{Criterion:} catch at least 70\% of eventual failures while flagging no more
than 20\% of the cohort, because staff can only call 20\%. \emph{Baseline:} flag
everyone who missed two or more labs.}

\lensPM{\emph{Decision:} whether to add a developer to a sprint at the mid-point.
\emph{Unit:} one sprint for one team. \emph{Target:} will committed story points
be delivered? \emph{Timing:} data up to day 5 of a 10-day sprint. \emph{Criterion:}
beat the scrum master's own mid-sprint judgement, measured over 20 sprints.
\emph{Baseline:} ``we finish if more than half the points are done by day 5.''}

\lensIOT{\emph{Decision:} whether to schedule a motor for maintenance this week.
\emph{Unit:} one motor on one day. \emph{Target:} will it fail within 30 days?
\emph{Timing:} only telemetry up to and including today. \emph{Criterion:} prevent
at least 80\% of unplanned failures with fewer than 5 unnecessary teardowns per
100 motors per year. \emph{Baseline:} fixed-interval servicing every 90 days.}

\lensSEC{\emph{Decision:} whether to force step-up authentication on a login.
\emph{Unit:} one authentication attempt. \emph{Target:} is this attempt an account
takeover? \emph{Timing:} only what is observable at the instant of the request.
\emph{Criterion:} detect 90\% of takeovers while challenging under 1 in 500
legitimate logins --- a false-positive budget set by the customer-experience team,
not by the data scientist. \emph{Baseline:} challenge every login from a new
country.}
\end{fourlenses}

%---------------------------------------------------------------------
\begin{lab}[Writing a Frame]
Take any dataset you have access to --- your own module grades will do --- and
fill in this template before writing any analysis code. It takes ten minutes and
saves days.

\begin{lstlisting}[language=Python]
FRAME = {
    "decision":  "What action changes because of this output?",
    "unit":      "One row is one ______.",
    "target":    "Exactly what is predicted/measured, in what units?",
    "timing":    "Which columns exist at prediction time? List them.",
    "criterion": "Success = ______ above/below ______.",
    "baseline":  "The no-model rule is ______, which achieves ______.",
}

for k, v in FRAME.items():
    print(f"{k.upper():10s} {v}")
    # Replace each value with your own answer, then re-read question 4.
\end{lstlisting}

\textbf{The test of a good frame:} hand it to someone outside the project. If they
can tell you what the model outputs and how you would know it worked, the frame is
done.
\end{lab}

%---------------------------------------------------------------------
\begin{checkpoint}
\begin{tightnum}
  \item A team reports that model accuracy rose from 94\% to 96\% on a dataset
        where 95\% of cases are negative. What should you ask before congratulating
        them?
  \item Why does CRISP-DM draw an arrow from deployment back to business
        understanding, rather than ending at deployment?
  \item A hospital wants to predict which patients will be readmitted within 30
        days. Name one feature that would be a leakage error, and say why.
\end{tightnum}
\end{checkpoint}

%---------------------------------------------------------------------
\begin{chaptersummary}
\begin{tight}
  \item CRISP-DM has six phases and is a loop; the back-steps are normal, not
        failures.
  \item Framing precedes data. Answer all six framing questions --- decision, unit,
        target, timing, criterion, baseline --- before opening a file.
  \item Timing (what is knowable at prediction time) prevents leakage, the most
        expensive routine mistake in the field.
  \item Model metrics and business metrics are different and can disagree. A
        project is judged on the latter.
  \item Modelling is the smallest phase. Framing and data preparation dominate.
\end{tight}
\end{chaptersummary}

\begin{reviewq}
  \item \textbf{(MCQ)} In CRISP-DM, which phase immediately follows data
        preparation? (a)~Evaluation (b)~Modelling (c)~Deployment (d)~Business
        understanding.
  \item \textbf{(MCQ)} A feature that would not be available when the model runs
        in production causes: (a)~overfitting (b)~leakage (c)~underfitting
        (d)~class imbalance.
  \item \textbf{(Short)} Give a definition of ``baseline'' and explain in one
        sentence why a project without one cannot demonstrate value.
  \item \textbf{(Short)} Describe a situation in which a model with higher accuracy
        delivers lower business value than a model with lower accuracy.
  \item \textbf{(Short)} Why is ``one row is one \emph{what}?'' a harder question
        than it appears for IoT telemetry?
  \item \textbf{(Applied)} Write a complete six-part frame for this request: ``We
        want to use AI to stop our servers going down.''
  \item \textbf{(Applied)} For the learning analytics lens above, propose a
        \emph{different} success criterion that would be appropriate if the
        department had unlimited staff, and explain how it changes which model you
        would prefer.
\end{reviewq}
